% Referencias verificadas el 10 de septiembre de 2026. No modifica datos.
\begin{thebibliography}{99}
\setlength{\itemsep}{0.10\baselineskip}
\phantomsection
\addcontentsline{toc}{section}{Referencias}

\bibitem{mustovicary}
B. Musto y J. Vicary,
``Quantum Latin squares and unitary error bases'',
\emph{Quantum Information and Computation} \textbf{16},
números 15--16, 1318--1332 (2016).
\href{https://doi.org/10.26421/QIC16.15-16-4}{doi:10.26421/QIC16.15-16-4}.

\bibitem{delascuevas2020}
G. De las Cuevas, T. Drescher y T. Netzer,
``Quantum magic squares: Dilations and their limitations'',
\emph{Journal of Mathematical Physics} \textbf{61}, 111704 (2020).
\href{https://doi.org/10.1063/5.0022344}{doi:10.1063/5.0022344}.
\href{https://arxiv.org/abs/1912.07332}{Prepublicación de los autores: arXiv:1912.07332}.

\bibitem{delascuevas2023}
G. De las Cuevas, T. Netzer e I. Valentiner-Branth,
``Magic squares: Latin, semiclassical, and quantum'',
\emph{Journal of Mathematical Physics} \textbf{64}, 022201 (2023).
\href{https://doi.org/10.1063/5.0127393}{doi:10.1063/5.0127393}.
\href{https://arxiv.org/abs/2209.10230}{Prepublicación de los autores: arXiv:2209.10230}.

\bibitem{codata2018}
E. Tiesinga, P. J. Mohr, D. B. Newell y B. N. Taylor,
``CODATA recommended values of the fundamental physical constants: 2018'',
\emph{Reviews of Modern Physics} \textbf{93}, 025010 (2021).
\href{https://doi.org/10.1103/RevModPhys.93.025010}{doi:10.1103/RevModPhys.93.025010}.
\href{https://physics.nist.gov/cuu/Constants/archive2018.html}{NIST: archivo del ajuste de 2018};
\href{https://physics.nist.gov/cuu/Constants/ArchiveASCII/allascii_2018.txt}{tabla ASCII archivada}.

\bibitem{codata2022}
P. J. Mohr, D. B. Newell, B. N. Taylor y E. Tiesinga,
``CODATA recommended values of the fundamental physical constants: 2022'',
\emph{Reviews of Modern Physics} \textbf{97}, 025002 (2025).
\href{https://doi.org/10.1103/RevModPhys.97.025002}{doi:10.1103/RevModPhys.97.025002}.
\href{https://physics.nist.gov/cuu/Constants/article2022.html}{NIST: documentación del ajuste de 2022}.
Las comparaciones del artículo conservan identificados los ajustes de 2018 y 2022.

\bibitem{pdg2026}
F. Takahashi et al. (Particle Data Group),
``Review of Particle Physics'',
\emph{International Journal of Modern Physics A} \textbf{41}, 2630011 (2026).
\href{https://doi.org/10.1142/S0217751X26300115}{doi:10.1142/S0217751X26300115}.
\href{https://pdg.lbl.gov/2026/}{Edición oficial de 2026};
\href{https://pdg.lbl.gov/2026/api/index.html}{documentación de acceso programático}.
El catálogo de este artículo conserva la instantánea identificada en las fuentes
como PDG 2026, con corte documental del 15 de enero de 2026 y consulta consignada
el 26 de julio de 2026. Sus 471 registros de masa y 384 de anchura, junto con la
huella de procedencia, se mantienen en el suplemento de datos; la referencia
web no sustituye esa instantánea.

\end{thebibliography}
